% NORMALIZACIONES DE TRANSCRIPCIÓN: los subíndices mecanografiados con letra o
% se representan como cero. La diagonal sin punta de (3.8.2) se conserva en la
% rama impresa; la rama actual restituye la flecha dirigida determinada por el
% texto.
% Canon: SGA-SRC-001262.
% El producto mecanografiado con X en (3.8.6) se representa mediante \times.
% La flecha de (3.8.10) sigue la autoridad francesa. No hay diferencias
% adicionales entre los testigos francés diplomático y corregido en este componente.

\subsection*{3.8. Biextensiones de $(P,Q)$ trivializadas sobre
extensiones de $P$ y $Q$}

Consideremos primero dos sucesiones exactas de grupos abelianos
\begin{equation}\tag{3.8.1}
 0\longrightarrow L_1\xrightarrow{\ i\ }L_0\xrightarrow{\ u\ }P
 \longrightarrow0,\qquad
 0\longrightarrow M_1\xrightarrow{\ j\ }M_0\xrightarrow{\ v\ }Q
 \longrightarrow0.
\end{equation}
Describiremos la categoría de las biextensiones $E$ de $(P,Q)$ por un
grupo abeliano $G$ cuya imagen inversa $(u,v)^*(E)$ es trivial; es
decir, aquellas para las que existe un homomorfismo $s$ que hace
conmutativo el triángulo siguiente:
\begin{equation}\tag{3.8.2}
\begin{tikzcd}[column sep=huge,row sep=huge]
 E\arrow[d]\sisi{\arrow[dr,dash,"s"]}{}&\\
 P\times Q&L_0\times M_0\arrow[l,"u\times v"']%
 \sisi{}{\arrow[ul,"s"]}
\end{tikzcd}\quad ,
\end{equation}
y que satisface la condición de aditividad requerida en cada uno de
sus dos argumentos. La elección de tal $s$ identifica, por tanto,
$(u,v)^*(E)$ con la biextensión trivial
$E_0=L_0\times M_0\times G$ de $(L_0,M_0)$ por $G$. La exactitud a
izquierda en $P,Q$ de la categoría cofibrada de las biextensiones de
$(P,Q)$ por $G$ fijo (3.7.6) muestra entonces que dar la biextensión
$E$ de $(P,Q)$, junto con la trivialización $s$ de su imagen inversa,
equivale esencialmente a dar dos trivializaciones parciales de la
biextensión trivial $E_0$ de $(L_0,M_0)$ por $G$, sobre $(L_1,M_0)$ y
$(L_0,M_1)$ respectivamente, que coincidan sobre $(L_1,M_1)$. Puesto
que las biextensiones inducidas por $E_0$ sobre $(L_1,M_0)$,
$(L_0,M_1)$ y $(L_1,M_1)$ se identifican con las biextensiones
triviales, una trivialización parcial sobre $(L_1,M_0)$
(respectivamente sobre $(L_0,M_1)$) equivale a dar un homomorfismo (o
emparejamiento)
\begin{equation}\tag{3.8.3}
 f:L_1\otimes M_0\longrightarrow G
 \qquad
 \bigl(\text{respectivamente }g:L_0\otimes M_1\longrightarrow G\bigr).
\end{equation}
La compatibilidad de que se trata significa simplemente que los dos
homomorfismos coinciden sobre $L_1\otimes M_1$, es decir,
\begin{equation}\tag{3.8.4}
 f(x_1\otimes y_1)=g(x_1\otimes y_1),
\end{equation}
para secciones $x_1,y_1$ de $L_1,M_1$ sobre un objeto de
$\underline T$. En términos del homomorfismo $s$ de (3.8.2), las
formas $f$ y $g$ quedan caracterizadas por
\begin{equation}\tag{3.8.5}
 s(x_1,y_0)=f(x_1\otimes y_0)e_{E/Q}(y_0),\qquad
 s(x_0,y_1)=g(x_0\otimes y_1)e_{E/P}(x_0),
\end{equation}
lo que pone de nuevo de manifiesto la necesidad de la condición
(3.8.4).

Elegida la trivialización $s$, las demás trivializaciones son los
morfismos $s':L_0\times M_0\longrightarrow E$ de la forma
\begin{equation}\tag{3.8.6}
 s'(x_0,y_0)=s(x_0,y_0)+h(x_0\otimes y_0),
\end{equation}
donde
\[
 h:L_0\otimes M_0\longrightarrow G
\]
es una forma arbitraria. De (3.8.5) se desprende entonces claramente
que las formas $f,g$ definidas por $s$ se transforman en
\begin{equation}\tag{3.8.7}
\begin{aligned}
 f'(x_1\otimes y_0)&=f(x_1\otimes y_0)+h(x_1\otimes y_0),\\
 g'(x_0\otimes y_1)&=g(x_0\otimes y_1)+h(x_0\otimes y_1).
\end{aligned}
\end{equation}
Así se obtiene:

\underline{Proposición} 3.8.8.
\underline{El subgrupo de} $\Biext^1(P,Q;G)$ \underline{formado por las clases
de biextensiones cuya imagen inversa por $(u,v)$ es una biextensión
trivial de $(L_0,M_0)$} (es decir, el núcleo de
$\Biext^1(P,Q;G)\longrightarrow\Biext^1(L_0,M_0;G)$)
\underline{es canónicamente isomorfo, mediante la descripción
precedente, al cociente del grupo de pares de formas $(f,g)$ de
(3.8.3) que satisfacen la condición de compatibilidad (3.8.4), por
los pares de formas que proceden por restricción de una forma}
$h:L_0\otimes M_0\longrightarrow G$.

Interpretemos ahora este isomorfismo en términos de álgebra homológica,
teniendo en cuenta el isomorfismo canónico
\[
 \Biext^1(P,Q;G)\simeq
 \Ext^1(P\mathbin{\overset{\mathbb L}{\otimes}}Q,G)
\]
de 3.6.5. Para ello, interpretemos los datos (3.8.1) como dos complejos
que resuelven $P$ y $Q$, respectivamente:
\begin{equation}\tag{3.8.9}
 L_\bullet:\ 0\longrightarrow L_1\longrightarrow L_0
 \longrightarrow0,\qquad
 M_\bullet:\ 0\longrightarrow M_1\longrightarrow M_0
 \longrightarrow0,
\end{equation}
y consideremos el complejo producto tensorial
$L_\bullet\otimes M_\bullet$
\[
 0\longrightarrow L_1\otimes M_1
 \xrightarrow{d_2}
 L_1\otimes M_0+L_0\otimes M_1
 \xrightarrow{d_1}
 L_0\otimes M_0\longrightarrow0.
\]
Se tiene, además, un homomorfismo canónico en la categoría derivada
\[
 P\mathbin{\overset{\mathbb L}{\otimes}}Q
 \longrightarrow L_\bullet\otimes M_\bullet,
\]
y, por tanto, un homomorfismo igualmente canónico
\begin{equation}\tag{3.8.10}
 H^1\!\left(\Hom^{\cdot}(L_\bullet\otimes M_\bullet,G)\right)
 \longrightarrow
 \Ext^1(P\mathbin{\overset{\mathbb L}{\otimes}}Q,G).
\end{equation}
El grupo de los $1$-cociclos de
$\Hom^{\cdot}(L_\bullet\otimes M_\bullet,G)$ se identifica con el
grupo de pares $(f,g)$ como en (3.8.3) que satisfacen la condición de
compatibilidad (3.8.4): a tal par se asocia el homomorfismo
$F:L_1\otimes M_0+L_0\otimes M_1\longrightarrow G$ de componentes
$(f,g)$,
\begin{equation}\tag{3.8.11}
 F(x_1\otimes y_0+x_0\otimes y_1)
 =f(x_1\otimes y_0)+g(x_0\otimes y_1).
\end{equation}
Además, el coborde de la $0$-cocadena
$h:L_0\otimes M_0\longrightarrow G$ es precisamente el par de sus
restricciones a $L_1\otimes M_0$ y $L_0\otimes M_1$. Por consiguiente,
el grupo de pares $(f,g)$ módulo los pares $(h,h)$ descritos en 3.8.8
es canónicamente isomorfo al primer miembro de (3.8.10). Sentado esto,
se tiene la compatibilidad siguiente.

\underline{Proposición} 3.8.12.
\underline{Si}
\[
 \xi\in\Ker\!\left(
 \Biext^1(P,Q;G)\longrightarrow\Biext^1(L_0,M_0;G)\right)
\]
\underline{está descrito por el par $(f,g)$ como en 3.8.8, entonces
la imagen de $\xi$ en}
$\Ext^1(P\mathbin{\overset{\mathbb L}{\otimes}}Q,G)$
\underline{mediante el isomorfismo (3.6.5.1) es la imagen por
(3.8.10) del elemento de su primer miembro definido por el cociclo
(3.8.11)}.

La comprobación de esta compatibilidad se deja al lector.

3.8.13. Aplicaremos esto, en particular, al caso
\[
 L_0=\mathbb Z[P],\qquad M_0=\mathbb Z[Q],
\]
en el que $L_1$ y $M_1$ son, respectivamente, los núcleos de los
homomorfismos canónicos $\mathbb Z[P]\longrightarrow P$ y
$\mathbb Z[Q]\longrightarrow Q$. Obsérvense los isomorfismos
canónicos
\[
 L_0\mathbin{\overset{\mathbb L}{\otimes}}M_0
 \simeq L_0\otimes M_0
 \simeq\mathbb Z[P\times Q].
\]
Como $L_0$ y $M_0$ son planos, (1.4.5) da isomorfismos canónicos
\[
 \Ext^i(L_0\mathbin{\overset{\mathbb L}{\otimes}}M_0,G)
 \simeq H^i(P\times Q,G_{P\times Q}).
\]
De aquí y de 3.6.5, o directamente de las definiciones y de 1.4, se
deduce que el funtor que asigna a cada biextensión de $(L_0,M_0)$ por
$G$ el $G_{P\times Q}$-torsor sobre $P\times Q$ que induce es una
equivalencia de categorías. (Esto sigue siendo válido sin utilizar
una estructura aditiva sobre $P,Q$.) Por consiguiente, si $E$ es una
biextensión de $(P,Q)$ por $G$, dar una trivialización de su imagen
inversa sobre $(L_0,M_0)$ equivale a dar una sección de $E$
considerado únicamente como objeto sobre $P\times Q$; es decir, una
trivialización de $E$ considerado como torsor bajo $G_{P\times Q}$.

Las biextensiones provistas de esta estructura adicional quedan, por
tanto, descritas por los pares $(f,g)$ de (3.8.3) que satisfacen la
compatibilidad (3.8.4), y su clasificación salvo isomorfismo es la
dada en 3.8.8.

Con la notación de 3.5, también puede considerarse $L_1$ como el
conúcleo de $L_2(P)\longrightarrow L_1(P)$, y análogamente $M_1$.
Esto permite expresar los datos $(f,g)$ mediante un sistema de
homomorfismos
\[
\left\{
\begin{aligned}
 f'&:L_1(P)\otimes L_0(Q)
     =\mathbb Z(P\times P\times Q)\longrightarrow G,\\
 g'&:L_0(P)\otimes L_1(Q)
     =\mathbb Z(P\times Q\times Q)\longrightarrow G,
\end{aligned}
\right.
\]
es decir, homomorfismos de haces de \underline{conjuntos}
\[
 P\times P\times Q\longrightarrow G,\qquad
 P\times Q\times Q\longrightarrow G,
\]
que satisfacen cinco condiciones de compatibilidad. Estas expresan
que $f'$ y $g'$ se factorizan a través de $f$ y $g$---lo que da
$2\times2$ ecuaciones, teniendo en cuenta la expresión de $L_2(P)$ y
$L_2(Q)$ como sumas de dos términos---junto con la condición (3.8.4).
Se ve entonces que $f'$ y $g'$ son precisamente los datos de
multiplicación parcial que definen una biextensión, y que las cinco
condiciones de que se trata son exactamente los cinco axiomas de las
biextensiones considerados en 2.1 (cf. la demostración de 3.6.4).

\subsection*{Bibliografía del Exposé VII}

\begin{enumerate}[label={[\arabic*]}]
\item M. Demazure, J. Giraud y M. Raynaud,
\emph{Schémas abéliens}, Séminaire Fac. Sc. Orsay 1967/68 (en prensa).
\item J. Giraud, \emph{Cohomologie non abélienne}, tesis, París,
1967 (en prensa).
\item A. Grothendieck, \emph{Sur quelques points d'Algèbre
homologique}, Tohoku Mathematical Journal 9 (\sisi{1956}{1957}), 119--221.
% Canon: SGA-SRC-001265.
\item A. Grothendieck, \emph{Catégories cofibrées additives et
Complexe cotangent relatif},
\sisi{Lectures Notes in Math.}{Lecture Notes in Mathematics} n.º 79, 1968
(Springer).
\item J.-L. Verdier, \emph{Catégories dérivées des catégories
abéliennes} (en prensa).
\item D. Mumford, \emph{Biextensions of formal groups}, Tata
Institute of Fundamental Research, 1968 (en prensa).
\end{enumerate}
